\newcommand*{\version}{0.2.0}

\documentclass[12pt,
              article,
              oneside
              ]{memoir}

\usepackage{iftex}

\ifluatex
  \usepackage{fontspec}
\fi

\usepackage[
  % latin,
  english,
  % spanish,% <- principal
  % es-minimal,
  % es-noshorthands,
  % es-tabla,
]{babel}

\ifluatex
% LuaLaTeX --------------------
\defaultfontfeatures{Ligatures=TeX, Scale=MatchLowercase}
\usepackage[MainFeatures={medium,
  Ligatures=Common,
  Numbers=OldStyle,
  SmallCapsFeatures = {LetterSpace=5.0}}]
{junicodevf}
%
% \setsansfont{Libertinus Sans}
% (sin las ligaduras de TeX: -{}-help no se vuelve una raya)
\setmonofont[Scale=MatchLowercase, Ligatures=ResetAll]{DejaVu Sans Mono}
% \setmonofont[Scale=MatchLowercase]{Libertinus Mono}
% \setmonofont[Scale=MatchLowercase}]{Fira Code}
\else
% PDFLaTeX ---------------------
\usepackage[T1]{fontenc}
\usepackage[osf]{libertinus}
\usepackage[scaled=0.9]{FiraMono}
% \usepackage[scaled=0.9]{cascadia-code}
% \usepackage[scaled=0.9]{roboto}
\input glyphtounicode
\pdfgentounicode=1
\fi

\usepackage[babel=true, verbose=false, tracking=true, expansion=true, protrusion=true, final, draft=false]{microtype}
\SetTracking{ encoding = *, shape = sc }{ 30 }
\SetTracking[context = notracking, ]{encoding = *}{0}

\usepackage[autostyle=false, style=british]{csquotes}
\usepackage{latexcolors}

\newcommand{\Rojo}{\color[HTML]{8B0000}}
\newcommand*{\rojo}[1]{\textcolor[HTML]{8B0000}{#1}}
\newcommand*{\rojoit}[1]{\textit{\textcolor[HTML]{8B0000}{#1}}}
\newcommand{\Azul}{\color{bluenivaca}}
\newcommand*{\azul}[1]{\textcolor{bluenivaca}{#1}}
\newcommand{\Verde}{\color{darkgreen}}
\newcommand*{\verde}[1]{\textcolor{darkgreen}{#1}}
\newcommand{\Amarillo}{\color{goldenpoppy}}
\newcommand*{\amarillo}[1]{\textcolor{goldenpoppy}{#1}}
\newcommand{\Morado}{\color{byzantium}}
\newcommand*{\morado}[1]{\textcolor{byzantium}{#1}}
\newcommand{\Naranja}{\color{orange}}
\newcommand*{\naranja}[1]{\textcolor{orange}{#1}}

% \usepackage[margin=3cm]{geometry}

\usepackage{graphicx}
\usepackage{booktabs}

\clubpenalty=9999
\widowpenalty=9999
\brokenpenalty=5000
\frenchspacing
\raggedbottom
\midsloppy
\parindent=14pt
\fboxsep=2pt

\renewcommand\thesection{\arabic{section}}
\chapterstyle{reparticle}

\usepackage{enumitem}
\setlist{noitemsep, nosep}
\setlist[description]{font=\itshape}

\newcommand*{\elipsis}{\,.\,.\,.\,}
\newcommand*{\elipsisb}{\textup{[\elipsis]}}

\renewenvironment{quote}%
 {\begin{list}{}%
              {\setlength\rightmargin{0pt}%
              \setlength\leftmargin{\parindent}}%
  \item[]\small\ignorespaces}
 {\unskip\end{list}}


\definecolor{Link}{HTML}{408080}
\usepackage{hyperref}
\hypersetup{
  colorlinks=true,
  linkcolor=blue,
  filecolor=blue,
  citecolor = black,
  urlcolor=Link,
}
\urlstyle{tt}
\usepackage{bookmark}

\setlength{\droptitle}{-2cm}

\title{\bfseries%
  \texttt{py-slowdown} v.\version{} \\
  User Manual}

\author{%
  Nicolás Vaughan\\
  \begin{minipage}[t]{0.8\linewidth}
    \centering\small%
    \href{mailto:nicolasvaughan@proton.me}{\texttt{nicolasvaughan@proton.me}}\\
    \url{https://nicolasvaughan.org}\\
    \url{https://github.com/nivaca/py-slowdown}\\
    \url{https://codeberg.org/nivaca/py-slowdown}\\
  \end{minipage}}


\date{2026-11-24}

% The pictures, made by tools/manual_images.py
\graphicspath{{img/manual/}}
\usepackage{longtable}

% the title and the author in the PDF's own metadata
\hypersetup{pdftitle={py-slowdown User Manual}, pdfauthor={Nicolás Vaughan}}

% the index (makeindex, with the style manual.ist: see latexmkrc)
\makeindex

% menus, keys, and the pictures of the window
\newcommand*{\menu}[1]{\textit{#1}}
\newcommand*{\then}{\,>\,}
\newcommand*{\key}[1]{\texttt{#1}}
\newcommand*{\shot}[3]{%
  \begin{figure}[htbp]
    \centering
    \includegraphics[width=#2\textwidth,height=0.75\textheight,keepaspectratio]{#1}
    \caption{#3}
    \label{fig:#1}
  \end{figure}}

\begin{document}
\maketitle

\tableofcontents*

\section{Introduction}
\noindent
\texttt{py-slowdown} plays a recording slower or faster without changing its
pitch. I made it for transcribing recorded lectures and interviews, where a
passage must be heard again and again at a speed one can follow; it serves
as well for dictations, for learning a language, and for learning music by
ear.

\index{waveform}\index{speed}\index{bookmarks}\index{formats}%
The recording is shown as a waveform, over which the position moves while
it plays. The usual controls move through it; the speed goes from 0.25\,×
to 3\,×, and can be changed while the recording plays; bookmarks, each with
a label, keep the places one wants to come back to. For each file the
program remembers its bookmarks and the point where it was left, and goes
on from there the next time.

\index{rubber band@Rubber Band}\index{quality}%
The stretching is done by the Rubber Band Library, with its R3 engine, the
one of highest quality, made for speech and music. At 1\,× the file is
played exactly as it is, sample for sample, without going through Rubber
Band.

\index{language of the interface}%
The window is in English or Spanish. This manual describes version
\version{}; the pictures are of the English window.

\section{Installing}
\label{sec:installing}

\index{installing}%
The releases are on GitHub and Codeberg (the addresses are on the first
page). Each one has the installers and packages below, and every file is
signed: see section~\ref{sec:signatures}. Each installer carries all the
program needs, its own Python and libraries included; nothing else has to
be installed.

\subsection{Windows}

\index{windows@Windows}\index{installer}%
\texttt{py-slowdown-\version-windows-x64-setup.exe} installs the program on
Windows 10 or later. It installs it for the user who runs it, with no
administrator's rights, unless one chooses to install it for everyone on
its first page. It adds an entry to the Start menu, an icon on the desktop
if one asks, and \menu{py-slowdown} in \menu{Open with} for audio files
(without taking from any other program the files it opens); the program is
removed in \menu{Settings\then Apps}, as any other.

\index{smartscreen@SmartScreen}%
The installer is not signed with a certificate that Windows knows, so the
first time Windows may say that it protected the computer
(SmartScreen). \menu{More info\then Run anyway} goes on. The signature of
section~\ref{sec:signatures} is what proves that the file is the one
published.

\subsection{macOS}

\index{macos@macOS}%
There is a disk image for each kind of Mac: \texttt{py-slowdown-\version-macos-arm64.dmg}
for those with Apple Silicon (M1 and later), and
\texttt{py-slowdown-\version-macos-x86\_64.dmg} for those with Intel. They
need macOS 13 (Ventura) or later. Open the image and drag
\menu{py-slowdown} to \menu{Applications}. Audio files can then be opened
with it from the Finder (\menu{Open With}), or dropped on its icon.

\index{gatekeeper@Gatekeeper}%
The application is not notarized by Apple, so the first time macOS says
that it cannot check it. Open \menu{System Settings\then Privacy \&
Security}, and at the bottom choose \menu{Open Anyway} for py-slowdown
(on macOS 14 and earlier, a right click on the application and
\menu{Open} does the same). It is asked only once.

\subsection{Any Linux: the installer}

\index{installer}%
\texttt{py-slowdown-\version-linux-x86\_64.run} installs the program on any
Linux with glibc 2.34 or later, with an entry in the menu of applications:

\begin{quote}
\texttt{sh py-slowdown-\version-linux-x86\_64.run}
\end{quote}

\noindent It installs it for the user who runs it (in
\texttt{\textasciitilde/.local/share/py-slowdown/app}, with the command
\texttt{py-slowdown} in \texttt{\textasciitilde/.local/bin}), or, run with
\texttt{sudo}, for everyone (in \texttt{/opt/py-slowdown}). It asks first
(\texttt{-{}-yes} does not ask), and \texttt{-{}-check} only checks the
file. To remove the program, run \texttt{uninstall.sh}, in its folder.

\subsection{Fedora}

\index{fedora (copr, rpm)@Fedora (Copr, RPM)}%
The simplest way, which also keeps the program up to date, is the Copr
repository:

\begin{quote}
\texttt{sudo dnf copr enable nivaca/py-slowdown}\\
\texttt{sudo dnf install py-slowdown}
\end{quote}

\noindent The RPM of each release can also be installed by itself
(\texttt{sudo dnf install ./py-slowdown-\version-1.fc45.noarch.rpm}). It
uses Fedora's Rubber Band, PortAudio and \texttt{ffmpeg}, which
\texttt{dnf} installs with it.

\subsection{Any Linux: the AppImage}

\index{appimage@AppImage}%
The AppImage is a single file that runs on any Linux with glibc 2.34 or
later (Ubuntu 22.04, Debian 12, Fedora 35 and later) and needs nothing
installed: it carries its own Python, Qt, FFmpeg, Rubber Band and
PortAudio. Make it executable and run it:

\begin{quote}
\texttt{chmod +x py-slowdown-\version-x86\_64.AppImage}\\
\texttt{./py-slowdown-\version-x86\_64.AppImage [FILE]}
\end{quote}

\subsection{Debian and Ubuntu}

\index{debian and ubuntu@Debian and Ubuntu}%
The Debian package, for Debian 12, Ubuntu 22.04 and later, is made from the
AppImage and installs it in \texttt{/opt/py-slowdown}, with the command
\texttt{py-slowdown}, an entry in the menu of applications and the manual
page:

\begin{quote}
\texttt{sudo apt install ./py-slowdown\_\version-1\_amd64.deb}
\end{quote}

\subsection{From the source, on Linux and macOS}

\index{installer}\index{macos@macOS}%
With the source (cloned, or the ZIP of a release unpacked) and Python 3.14
or later, \texttt{install/install.sh} installs the program for one user, in
an environment of its own, with an entry in the menu of applications (on
macOS, an application in \texttt{\textasciitilde/Applications}); run again,
it updates it, and \texttt{install/install.sh -{}-uninstall} removes it. The
program then needs, from the system, the Rubber Band Library (version 3 or
later) and, on Linux, PortAudio:

\begin{quote}
Fedora: \texttt{sudo dnf install rubberband-libs portaudio}\\
Debian, Ubuntu: \texttt{sudo apt install librubberband3 libportaudio2}\\
(\texttt{librubberband2} on older releases)\\
macOS: \texttt{brew install rubberband}
\end{quote}

\noindent\index{uv}%
To run it from the source without installing it, with \texttt{uv}:
\texttt{uv run py-slowdown [FILE]}.

\subsection{Checking the files of a release}
\label{sec:signatures}

\index{signatures}%
Each file of a release has a signature (\texttt{FILE.asc}), and
\texttt{SHA256SUMS} lists the sums of all of them, with its own signature.
They are checked with the public key published with the release:

\begin{quote}
\texttt{gpg -{}-import RPM-GPG-KEY-py-slowdown}\\
\texttt{gpg -{}-verify FILE.asc FILE}
\end{quote}

\noindent The installers for Windows and macOS, which are made by GitHub,
have their sums in \texttt{SHA256SUMS-installers}, also signed. The RPM is
signed inside too; \texttt{dnf} checks it once the key is imported
(\texttt{sudo rpm -{}-import RPM-GPG-KEY-py-slowdown}).

\section{Starting the program}

\index{starting}\index{command line}%
The program is in the menu of applications, as \menu{py-slowdown}. From a
terminal, \texttt{py-slowdown} opens the window, and
\texttt{py-slowdown FILE} opens it with that file;
\texttt{py-slowdown -{}-help} prints the help, and \texttt{-{}-version} the version.
An audio file can also be dropped on the window, or opened with
\menu{File\then Open} (\key{Ctrl+O}) or \menu{File\then Open recent}.

\section{The window}

\shot{window_en}{1}{The window, with a recording at 0.75\,× and two bookmarks}

\index{window}\index{waveform}\index{controls}\index{bookmarks!panel}%
The waveform takes most of the window (figure~\ref{fig:window_en}). Above
it is a ruler of times; the red line is the position, the point that is
heard; what was already played is darker. The bookmarks are dashed lines,
with their labels on top. Under the waveform are the controls and the
time, and under them the speed, the pitch and the volume; on the right, the panel of the bookmarks, which
can be closed, moved or left floating (\menu{View\then Bookmarks} brings it
back). The status bar says what is being read, and then the format of the
file: its codec, its sample rate, mono or stereo, and its length.

\section{Opening a file}

\index{formats}\index{video files}%
The program opens whatever FFmpeg reads: MP3, M4A and AAC, MP4, FLAC, WAV,
AIFF, Ogg, Opus, WMA and many more, and the sound of video files (MP4, MKV,
WebM, MOV). A file with more than two channels is mixed down to stereo.

\index{long files}\index{reading}%
A file is read in the background. While it is read, the status bar shows
how much is done, the waveform fills in from the left, and what is already
read can be played. A lecture of two hours takes a few seconds. The
decoded sound is kept in a temporary file in the cache folder, not in the
memory (two hours of stereo take about 2.5\,GB there), and the temporary
file goes away when the file is closed, or when the program ends.

\index{position, remembered}%
When a file is closed (or another one is opened, or the program ends), the
program keeps the point where it was, and the next time the file is opened
it goes on from there. A file left at its end starts again from the
beginning.

\section{Playing}

\shot{controls_en}{1}{The controls}

\index{controls}\index{playing}\index{keys}%
The controls (figure~\ref{fig:controls_en}) are, from the left:

\begin{description}
\item[Go to the start] (\key{Home}).
\item[Back] by the skip length (\key{←}); held down, it keeps going back,
  to run through the file.
\item[Play and pause] (\key{Space}).
\item[Stop] and go back to the start (\key{S}).
\item[Forward] by the skip length (\key{→}); held down, it keeps going.
\item[Go to the end] (\key{End}).
\item[Bookmark] adds one at the position (\key{B}).
\item[Repeat] the selection, again and again (\key{L}; section~\ref{sec:repeat}).
\item[Count-in] before playing (\key{C}; section~\ref{sec:countin}).
\end{description}

\noindent\index{skip length}%
The skip length is 5 seconds unless another is chosen in
\menu{Settings\then Skip length} (1 to 60 seconds).
\key{Shift+←} and \key{Shift+→} go back and forward one second, which is
what is most often needed when transcribing. At the end of the file the
playing stops; \key{Space} then plays it again from the start.

\index{volume}%
The volume is the slider on the right. It only lowers the sound of the
program, never raises it above that of the file.

\section{The speed}

\index{speed}%
The speed goes from 0.25\,× (four times slower) to 3\,× (three times
faster). It is changed with the slider, by writing it in the box beside it,
or with the keys \key{↓} and \key{↑} (or \key{−} and \key{+}), by
steps of 0.05; \key{N}, or the button \menu{1×}, brings back the normal
speed. It can be changed while the recording plays; the change is heard
after about 0.15 seconds, the sound already prepared at the old speed.

\index{quality}\index{rubber band@Rubber Band}%
The pitch does not change: a voice at 0.5\,× says the words more slowly,
but with its own voice. The quality is that of Rubber Band's R3 engine,
which keeps speech clear and music natural, even far from 1\,×; at very
slow speeds the attacks of the sounds (consonants, the strokes of a
guitar) are somewhat softened, as with any method of its kind. At 1\,× the
file is played as it is.

\index{speed!not available}%
If the Rubber Band Library is not installed (which can happen only when the
program is installed from the source), the window says so in the status
bar, and the speed stays at 1\,×.

\section{The pitch}
\label{sec:pitch}

\index{pitch}\index{semitones|see{pitch}}\index{transposing|see{pitch}}%
The pitch can be moved up or down, in semitones, without changing the
speed: to play along with a recording in another key, to bring a song
within reach of a voice, or to hear a fast passage slower and an octave
higher, where its notes stand out. The box \menu{Pitch}, beside the speed,
goes from 12 semitones down to 12 up (an octave each way), by steps of
one; \key{Ctrl+↑} and \key{Ctrl+↓} do the same, and \key{Shift+N}, or
the button \menu{0}, brings back the pitch of the recording. The box also
takes tenths, written in it (0.5, say): a recording that is a little off
pitch, as old tapes often are, can be tuned to the instrument.

\index{quality}\index{rubber band@Rubber Band}%
The pitch can be changed while the recording plays, without a click. It
is done by Rubber Band too, in the same pass as the speed, with the same
quality: measured, a tone lands within a cent (a hundredth of a semitone)
of the pitch asked for, also at the slowest and fastest speeds. At the pitch of the
recording and at 1\,× the file is played as it is. The pitch, like the
speed, is not kept from one session to the next: the program always
starts at the recording's own.

\section{The waveform}

\index{waveform}\index{position}%
A click on the waveform takes the position there, and dragging moves it,
whether the recording plays or not. The waveform is scaled to the loudest
point of the file, so that a quiet recording does not look flat.

\shot{zoom_en}{0.85}{A closer view, around a bookmark}

\index{zoom}\index{follow the position}%
\key{Ctrl+wheel} zooms in or out around the pointer
(figure~\ref{fig:zoom_en}); \key{Ctrl++} and \key{Ctrl+−} do the same
around the position, and \key{Ctrl+0} shows the whole file again. Up close
the samples themselves are drawn, down to one per pixel. When the view is
closer than the whole file, a scroll bar appears under it, and the wheel
moves the view; while the recording plays, the view turns to the next page
when the position leaves it, unless \menu{View\then Follow the position}
is turned off.

\section{Bookmarks}

\index{bookmarks}%
A bookmark marks a point of the file, with a label if one wants. \key{B},
the button \menu{Bookmark} or \menu{Bookmarks\then Add a bookmark} put one
at the position; a right click on the waveform puts one where the pointer
is (figure~\ref{fig:context_en}). Two bookmarks are never put at the same
point: if there is one already, that one is chosen.

\shot{context_en}{0.35}{The menu of a right click over a bookmark}

\index{bookmarks!label}%
The label is written in the panel, with a double click on it, or in the
dialog that a double click on the bookmark in the waveform opens, as does
\menu{Rename the bookmark} in the menu of the right click
(figure~\ref{fig:rename_en}). A bookmark without a label shows its number.

\shot{rename_en}{0.35}{The label of a bookmark}

\shot{panel_en}{0.45}{The panel of the bookmarks}

\index{bookmarks!panel}\index{bookmarks!going to}%
The panel (figure~\ref{fig:panel_en}) lists the bookmarks in order, with
their times. A double click on a time takes the position there; a click on
a bookmark in the waveform, on its label, does the same.
\key{Ctrl+←} and \key{Ctrl+→} (or \key{PgUp} and \key{PgDn}) go to the
previous and the next bookmark. While the recording plays, the previous one
is the one before the bookmark just passed, so that a passage can be heard
again from its bookmark with one key.

\index{bookmarks!removing}\index{bookmarks!exporting}%
\menu{Remove} in the panel, or \menu{Remove the bookmark} in the menu of
the right click, removes one. \menu{Export} in the panel, or
\menu{Bookmarks\then Export the bookmarks}, saves them as a text file, one
per line: the time, a tab, and the label. That text can be pasted into a
transcription or a spreadsheet.

\index{bookmarks!kept}%
The bookmarks are saved as soon as they change, for each file, by its full
path (section~\ref{sec:files-kept}). The audio file is never changed. A
file that is moved or renamed is a new file for the program, without
bookmarks.

\section{Repeating a selection}
\label{sec:repeat}

\shot{repeat_en}{0.85}{A selection repeated, with the count-in at its third click}

\index{selection}\index{repeat}\index{loop|see{repeat}}%
A passage can be repeated without end, to hear it until it is understood
or to play along with it. First select it: \key{Shift+drag} on the
waveform; or, with the position at its start, \key{I}, and at its end,
\key{O} (one of them alone selects from the position to the end, or from
the start, of the file); or, with a right click, \menu{Select from here to
the next bookmark}. The selection is drawn in green, also on the ruler,
and the status bar says where it begins and ends. Then turn on
\menu{Playback\then Repeat the selection} (\key{L}), or the button with
the two arrows: when the playing reaches the end of the selection, it goes
on from its start (figure~\ref{fig:repeat_en}).

\index{repeat!without a gap}%
There is no pause and no click between the end and the start: the last
hundredth of a second of the selection fades into what comes before its
start, so that the passage sounds as one continuous stream. This also
holds at any speed, because the stretching is done on that stream, not on
each repetition apart. A click outside the selection takes the position
there as usual; if it is before the selection, the repeating begins when
the playing reaches it. \key{Esc}, or \menu{Clear the selection}, removes
it; \key{L} again stops the repeating and keeps the selection.

\section{Counting in}
\label{sec:countin}

\shot{countin_en}{0.45}{The settings of the count-in}

\index{count-in}\index{metronome|see{count-in}}\index{tempo}%
Before the recording starts, the program can count, as a metronome does
before a band comes in: a few clicks at a tempo, the first one higher, and
then the recording, on the next beat. Turn it on with
\menu{Playback\then Count in before playing} (\key{C}), or the button
\menu{Count-in}; it then counts each time playing starts (with
\key{Space} or the button), not when the position moves while it plays.
While it counts, the waveform shows the beat in large figures
(figure~\ref{fig:repeat_en}), and the position does not move.

\index{count-in!offset}\index{offset of the recording}%
The recording comes in on the beat after the last click if its first beat
is exactly at the position. When it is not (a bookmark placed a little
before the first beat, say), \menu{Recording offset}, in the same dialog,
moves it: a negative value brings it in earlier, a positive one later, by
that many seconds of the recording (from 2 earlier to 2 later, by
hundredths). If the first beat is 0.2 seconds after the position, $-0.2$
puts it on the beat, at any speed. It can come in earlier only up to the
last click, which is never cut: the dialog says when more is asked for.

The offset is easier to set on the waveform: with the position where
playing will start, a right click on the recording's first beat, and
\menu{Count-in: the first beat is here}. The program sets the offset from
the position to that point, turns the count-in on, and says what it did; a
green dashed line marked 1 then shows, while it does not play, where the
count-in expects the first beat (figure~\ref{fig:zoom_en}), to check it or
to set it again closer in. \menu{Count-in: the first beat at the position}
takes the offset back to nothing.

\index{tap tempo}\index{tempo!tapped}%
The tempo of a recording can be measured by tapping: while it plays,
press \menu{Tap}, beside \menu{Count-in}, or the key \key{T}, with the
beat, a few times. From the second tap on, the status bar says the tempo
measured, and it becomes that of the count-in; the more taps, the steadier,
since it is the middle value of the intervals between the last eight, so
that one uneven tap does not move it. A pause of more than two seconds
begins a new measure. What is tapped is the tempo heard: with
\menu{Follow the speed}, a beat tapped at 60 BPM while the recording
plays at 0.5\,× is kept as 120 BPM, the recording's own, and the
count-in is then right at any speed. The dialog of the count-in has a
\menu{Tap} button too, beside the tempo.

\index{bluetooth@Bluetooth}%
The first click comes four tenths of a second after Play, and the count
begins with a very faint hiss, hardly audible, that goes on between the
clicks: many Bluetooth earphones and headphones mute themselves while they
receive silence, and take a moment to wake, so that, without it, they
swallowed the short clicks of the count.

The arrow beside the button \menu{Count-in}, \menu{Settings\then
Count-in}, or \menu{Playback\then Count-in settings} open its settings
(a right click on the button too), and the button's tooltip says what
they are now. They set how many counts (from 1 to 16) and the tempo, in beats per minute
(from 20 to 300; figure~\ref{fig:countin_en}). With \menu{Follow the
speed}, the tempo changes with the speed: a song at 120 BPM played at
0.5\,× is counted at 60, as it will sound. The dialog says what will be
heard at the speed of that moment. These settings are kept.

\section{Settings}

\shot{menu_settings_en}{0.35}{The menu Settings}

\index{settings}\index{language of the interface}\index{skip length}%
The menu \menu{Settings} (figure~\ref{fig:menu_settings_en}) has the
settings of the count-in (section~\ref{sec:countin}), the skip
length of \menu{Back} and \menu{Forward}, and the language of the
interface. The language follows that of the system unless one is chosen;
a new one is used when the program starts again. The program also
remembers the size and place of the window and of its panel, the volume,
whether the view follows the position, the folder of the last file opened,
and the recent files.

\section{The files it keeps}
\label{sec:files-kept}

\index{configuration files}\index{bookmarks!kept}%
\begin{description}
\item[The settings] in
  \texttt{\textasciitilde/.config/py-slowdown/py-slowdown.conf} on Linux;
  in the registry, under \texttt{HKEY\_CURRENT\_USER\textbackslash
  Software\textbackslash py-slowdown}, on Windows; and in
  \texttt{\textasciitilde/Library/Preferences/com.py-slowdown.py-slowdown.plist}
  on macOS. Deleting them, with the program closed, brings back the
  defaults.
\item[The bookmarks] and the point where each file was left, in
  \texttt{bookmarks.json}: in \texttt{\textasciitilde/.local/share/py-slowdown/}
  on Linux (in \texttt{\$XDG\_DATA\_HOME}, if it is set), in
  \texttt{\%APPDATA\%\textbackslash py-slowdown} on Windows, and in
  \texttt{\textasciitilde/Library/Application Support/py-slowdown/} on
  macOS. It is a JSON file that can be copied to another computer, or read
  by other programs.
\item[The decoded sound] of the file open, in a temporary file without a
  name, which goes away with it: in \texttt{\textasciitilde/.cache/py-slowdown/}
  on Linux, \texttt{\%LOCALAPPDATA\%\textbackslash py-slowdown\textbackslash
  cache} on Windows, and \texttt{\textasciitilde/Library/Caches/py-slowdown/}
  on macOS.
\end{description}

\noindent Uninstalling the program never removes them.

\section{Keys}

\index{keys}%
\begin{longtable}{@{}p{0.42\textwidth}p{0.52\textwidth}@{}}
\toprule
\textbf{Keys} & \textbf{What they do}\\
\midrule
\endhead
\key{Space} & play or pause\\
\key{S} & stop and go back to the start\\
\key{Home}, \key{End} & go to the start, to the end\\
\key{←}, \key{→} & back, forward (the skip length)\\
\key{Shift+←}, \key{Shift+→} & back, forward one second\\
\key{↓}, \key{↑} (or \key{−}, \key{+}) & slower, faster (by 0.05\,×)\\
\key{N} & normal speed\\
\key{Ctrl+↑}, \key{Ctrl+↓} & a semitone higher, lower\\
\key{Shift+N} & the original pitch\\
\key{B} (or \key{Ctrl+B}) & add a bookmark at the position\\
\key{Ctrl+←}, \key{Ctrl+→} (or \key{PgUp}, \key{PgDn}) & previous, next bookmark\\
\key{Ctrl++}, \key{Ctrl+−}, \key{Ctrl+0} & zoom in, zoom out, the whole file\\
\key{Ctrl+wheel} & zoom around the pointer\\
\key{L} & repeat the selection\\
\key{Shift+drag} & select\\
\key{I}, \key{O} & start, end of the selection at the position\\
\key{Esc} & clear the selection\\
\key{C} & count in before playing\\
\key{T} & tap the tempo, with the beat\\
\key{Ctrl+O} & open a file\\
\key{Ctrl+W} & close the file\\
\key{F1} & the list of the keys\\
\bottomrule
\end{longtable}

\section{When something goes wrong}

\index{troubleshooting}%
\begin{description}
\item[There is no sound] check the volume of the program and that of the
  system, and that the right output is chosen in the settings of the
  system's sound; the program plays through the default output.
\item[The speed or the pitch cannot be changed] the Rubber Band Library is not
  installed (the status bar says so): install it as in
  section~\ref{sec:installing}.
\item[A file is not opened] FFmpeg could not read it: it may be damaged,
  or not be audio. A file damaged halfway is opened up to where it can be
  read, and the status bar says so.
\item[The bookmarks of a file are gone] the file was moved or renamed; its
  bookmarks are still in \texttt{bookmarks.json}, under its old path, and
  can be moved to the new one there, with the program closed.
\item[A setting is wrong] delete the settings
  (section~\ref{sec:files-kept}).
\item[macOS or Windows will not open the program] it is not signed by
  Apple or Microsoft: see section~\ref{sec:installing}.
\end{description}

\noindent
Problems and wishes can be written as issues on GitHub or
Codeberg (the addresses are on the first page), or sent to
\href{mailto:nicolasvaughan@proton.me}{\texttt{nicolasvaughan@proton.me}}.

\printindex

\end{document}

% ===============================================
%%% Local Variables:
%%% mode: latex
%%% TeX-master: t
%%% End:
